\section{总结与延伸}

\subsection{核心要点回顾}

本节课从推理的工作负载分析出发，系统性地介绍了优化推理效率的多种技术：

\begin{importantbox}{推理优化技术全景}
\begin{enumerate}
    \item \textbf{理解工作负载}：推理（特别是 Generation 阶段）是内存受限的，算术强度低
    \item \textbf{有损优化——缩减 KV Cache}：GQA、MLA、CLA、局部注意力
    \item \textbf{有损优化——替代架构}：SSM（Mamba 等）、扩散模型
    \item \textbf{有损优化——量化与剪枝}：int8/int4 量化、AWQ、模型剪枝+蒸馏
    \item \textbf{无损加速}：投机采样（保证精确采样）
    \item \textbf{系统优化}：连续批处理、PagedAttention
\end{enumerate}
\end{importantbox}

\subsection{关键洞察}

\begin{enumerate}
    \item \textbf{推理不仅仅是``运行模型''}：推理优化本质上是在给定资源预算下最大化输出质量。这意味着架构设计、模型训练和推理部署需要协同考虑。

    \item \textbf{最大的优化潜力在于架构创新}：系统级优化（更好的 kernel、更高效的调度）能带来常数级的改进。而架构创新（SSM、扩散模型）可以改变计算复杂度的\textbf{渐近行为}，潜力远大于系统优化。

    \item \textbf{借鉴经典计算机科学}：投机采样借鉴了处理器中的\textbf{投机执行}（speculative execution），PagedAttention 借鉴了操作系统的\textbf{虚拟内存分页}（paging）。
\end{enumerate}

\subsection{拓展阅读}

\begin{itemize}
    \item \textbf{Scaling Book --- Transformers 章节}：\url{https://jax-ml.github.io/scaling-book/transformers/}
    \item \textbf{Scaling Book --- Inference 章节}：\url{https://jax-ml.github.io/scaling-book/inference/}
    \item \textbf{GQA 论文}：Ainslie et al., ``GQA: Training Generalized Multi-Query Transformer Models from Multi-Head Checkpoints''
    \item \textbf{MLA / DeepSeek v2}：DeepSeek-AI, ``DeepSeek-V2: A Strong, Economical, and Efficient Mixture-of-Experts Language Model''
    \item \textbf{Mamba}：Gu and Dao, ``Mamba: Linear-Time Sequence Modeling with Selective State Spaces'' (\url{https://arxiv.org/abs/2312.00752})
    \item \textbf{投机采样}：Leviathan et al. (\url{https://arxiv.org/abs/2211.17192}); Chen et al. (\url{https://arxiv.org/abs/2302.01318})
    \item \textbf{Medusa}：\url{https://arxiv.org/abs/2401.10774}
    \item \textbf{EAGLE}：\url{https://arxiv.org/abs/2401.15077}
    \item \textbf{LLM.int8()}：Dettmers et al. (\url{https://arxiv.org/abs/2208.07339})
    \item \textbf{AWQ}：Lin et al. (\url{https://arxiv.org/abs/2306.00978})
    \item \textbf{vLLM / PagedAttention}：Kwon et al. (\url{https://arxiv.org/abs/2309.06180})
    \item \textbf{Orca (Continuous Batching)}：Yu et al. (\url{https://www.usenix.org/system/files/osdi22-yu.pdf})
    \item \textbf{character.ai 推理优化}：\url{https://research.character.ai/optimizing-inference/}
\end{itemize}

\end{document}
